\section{Introducción: por qué el tono de este informe trimestral es ambivalente}

El tono central de este episodio del «informe trimestral global de modelos grandes» tiene dos caras. Por un lado, Coding y Agentic workflow están llevando a la IA de los chatbots a sistemas capaces de realizar trabajo, y tanto la capacidad de los modelos como la eficiencia de la investigación se aceleran. Por otro lado, esta aceleración está empujando el trabajo de oficina, la organización de la investigación y desarrollo y la distribución social hacia un periodo de deflación y desempleo. El juicio central de Guangmi (广密) es que Coding es el nuevo acelerador de la IA: un modelo líder en Coding es como una GPU de vanguardia y se convertirá en un amplificador decisivo del avance hacia la AGI.

\begin{importantbox}{Tesis central de este episodio}
Si el Chatbot es el primer acto y Coding/Agent es el segundo, en la siguiente etapa las empresas de modelos no compiten solo por «quién conversa mejor», sino por quién logra conectar el modelo con el código, las herramientas, los entornos de tareas y los flujos de trabajo reales, y convertir esa capacidad en una nueva generación de sistema operativo.
\end{importantbox}

\begin{knowledgebox}{Nota sobre la estrategia visual}
Este video muestra la imagen fija de una entrevista o pódcast de audio, sin slides didácticas, pizarrón ni demostraciones de producto. Conforme al estándar de pódcast de este repositorio, el cuerpo del texto no repite fotogramas de los participantes; la portada sirve para identificar la fuente y el cuerpo del texto expone la estructura del panorama mediante diagramas conceptuales y tablas.
\end{knowledgebox}

\subsection{Resumen de la sección}

Este episodio no es una crónica de noticias, sino una actualización trimestral del marco de análisis: Coding se convierte en acelerador de la AGI, se reordenan los periodos de competencia entre las empresas de modelos, el modelo podría convertirse en la nueva generación de OS y el impacto social empieza a pasar de la «eficiencia de las aplicaciones» a la «estructura de los puestos de trabajo».
